% =====================================================================
% EXERGUE ET REMERCIEMENTS
%
% Fichier PARTAGÉ par les quatre versions du document, au même titre que les
% chapitres et que la page de garde : main.tex, main_v2.tex, main_v3.tex et
% main_ensae.tex l'appellent tous. NE PAS LE DUPLIQUER pour faire évoluer une
% version, ce serait quatre textes à maintenir.
%
% UN CHAMP RESTE À COMPLÉTER, ET IL N'APPARTIENT PAS À L'ASSISTANT : le nom de
% famille du fondateur de Nexialog Consulting. Il s'imprime en gras entre
% crochets tant qu'il n'est pas renseigné, donc un dépôt avec un nom absent est
% visible à la relecture plutôt que silencieux.
% =====================================================================

% --- Le champ à compléter --------------------------------------------
\newcommand{\fondateurnexialog}{Ali~BEHBAHANI}
% ---------------------------------------------------------------------

% =====================================================================
%  EXERGUE
%
%  La citation est la « loi de la crédibilité décroissante » de Manski, qui
%  ouvre Identification Problems in the Social Sciences. Elle est à sa place
%  ici et pas seulement par ornement : le mémoire refuse de poser la direction
%  de la contagion et publie un capital en bornes, ce qui est exactement
%  l'arbitrage qu'elle énonce.
%
%  ELLE EST CITÉE DANS SA LANGUE D'ORIGINE DEPUIS LE 28 SEPTEMBRE 2026, demande
%  de Kélian. Elle portait jusque-là une traduction de l'auteur, signalée comme
%  telle ; l'anglais évite d'avoir à défendre une traduction qui n'apporte rien,
%  la phrase étant courte et le lecteur d'un mémoire d'actuariat la lisant sans
%  peine. Pour mémoire, la traduction retirée : « La crédibilité d'une inférence
%  décroît avec la force des hypothèses que l'on maintient. »
%
%  LA CÉSURE EST INTERDITE DANS CE BLOC, et ce n'est pas cosmétique : babel est
%  chargé en français SEUL, donc il couperait cette phrase anglaise sur des
%  motifs français. Interdire la césure ici coûte deux lignes bien réparties et
%  évite de toucher au préambule à deux jours du dépôt.
% =====================================================================
\thispagestyle{empty}
\vspace*{0.30\textheight}

\begin{flushright}
\begin{minipage}{0.72\textwidth}
\itshape\hyphenpenalty=10000 \exhyphenpenalty=10000
The credibility of inference decreases with the strength of the assumptions
maintained.

\vspace{10pt}
\normalfont\small\raggedleft
Charles F. Manski, \emph{Identification Problems in the Social Sciences}, 1995.
\end{minipage}
\end{flushright}

\clearpage

% =====================================================================
%  REMERCIEMENTS
% =====================================================================
\chapter*{Remerciements}
\addcontentsline{toc}{chapter}{Remerciements}
% Voir la note de 20_notes_synthese : un \chapter* laisse le titre courant sur
% la marque precedente. Sans effet visible ici, la page etant unique, mais la
% regle vaut pour les quatre versions et le texte peut s'allonger.
\markboth{Remerciements}{}

L'aboutissement de ce mémoire doit beaucoup à des relectures exigeantes, et à des
personnes qui ont accepté de discuter ce qu'il affirmait plutôt que la façon dont il le
calculait. Je souhaite leur exprimer ici ma gratitude.

Je tiens en premier lieu à remercier \fondateurnexialog, fondateur de Nexialog
Consulting, de m'avoir accueilli au sein de son cabinet et pour la confiance qu'il m'a
témoignée. Un travail de cette nature demande du temps long, et peu d'endroits
l'accordent à un stagiaire.

Je souhaite remercier tout particulièrement mon maître de stage, Hugo RAPIOR,
responsable R\&D en actuariat, pour son accueil au sein de son équipe, sa disponibilité
et la qualité de son encadrement. Nos points d'avancement hebdomadaires et ses conseils
ont beaucoup apporté à ce travail. Ces quelques mois m'auront appris autant sur la
conduite d'un travail que sur son contenu.

J'adresse également mes remerciements à l'ensemble des équipes du pôle Actuariat de
Nexialog Consulting, pour leur accueil et la qualité de nos échanges, et plus
spécialement aux trois praticiens qui ont pris le temps de confronter les affirmations
du modèle à leur expérience.

% LES TROIS RELECTEURS NE SONT PAS NOMMÉS, et cette règle ne bouge pas : décision de
% Kélian du 9 septembre 2026, tenue alors même que la première répondante avait coché
% « citez-moi par mon nom ». Ils ne sont désignés ici que par leur nombre, et par leur
% rôle au § de la relecture de praticien de l'annexe C, où leurs réponses sont rapportées.
% Le 28 septembre, la phrase ne citait que la première : elle les cite tous les trois.

Mes remerciements vont aussi au corps enseignant de l'ENSAE Paris et à la filière
Actuariat, pour la qualité de leurs enseignements. Je remercie en particulier Caroline
HILLAIRET, ma tutrice académique, pour son suivi et pour ses remarques, notamment sur le
traitement de l'incertitude, qui ont fait évoluer ce mémoire.

À ma famille enfin, je veux dire ma reconnaissance la plus sincère. Son soutien n'a
jamais eu besoin d'être demandé ni expliqué, et c'est la forme de soutien qui compte le
plus. Ce travail lui revient autant qu'à moi.

\clearpage
